\documentclass[12pt]{article}
\usepackage[utf8]{inputenc}
\usepackage[T1]{fontenc}
\usepackage{amsmath, amssymb}
\usepackage{graphicx}
\usepackage{hyperref}
\usepackage{geometry}
\usepackage{listings}
\usepackage{xcolor}
\geometry{margin=1in}

\lstset{
    basicstyle=\ttfamily\small,
    breaklines=true,
    frame=single,
    backgroundcolor=\color{gray!10},
    keywordstyle=\color{blue},
    commentstyle=\color{gray},
    stringstyle=\color{teal}
}

\title{Web Technologies 1 (Frontend) \\ Week 2: HTML Advanced}
\author{Dossymzhan}
\date{}

\begin{document}
\maketitle

\section{Overview}
Week 2 goes beyond the basic HTML skeleton: tables, forms/inputs/buttons, lists/links/images, semantic tags, \texttt{<div>} and \texttt{<span>}, and finally how to publish a site with GitHub Pages.

\section{1. Tables}
Tables organize related information into rows and columns.
\begin{lstlisting}[language=HTML]
<table>
  <thead>
    <tr><th>Subject</th><th>Day</th></tr>
  </thead>
  <tbody>
    <tr><td>Web Tech 1</td><td>Monday</td></tr>
    <tr><td>Algorithms</td><td>Tuesday</td></tr>
  </tbody>
  <tfoot>
    <tr><td colspan="2">Total: 2 classes</td></tr>
  </tfoot>
</table>
\end{lstlisting}
\begin{itemize}
    \item \texttt{<tr>} — table row
    \item \texttt{<th>} — table header cell
    \item \texttt{<td>} — table data cell
    \item \texttt{colspan} lets a cell span multiple columns.
\end{itemize}

\section{2. Forms, Inputs and Buttons}
The \texttt{<form>} element collects user input.
\begin{lstlisting}[language=HTML]
<form>
  <label>Name</label>
  <input type="text">

  <label>Email</label>
  <input type="email">

  <button>Send</button>
</form>
\end{lstlisting}
\begin{itemize}
    \item An \texttt{<input>} element's appearance/behavior depends on its \texttt{type} attribute.
    \item Common input types: \texttt{text}, \texttt{email}, \texttt{password}, \texttt{number}, \texttt{date}, \texttt{color}, \texttt{checkbox}, \texttt{radio}, \texttt{file}, \texttt{submit}.
    \item \texttt{<label>} defines a caption for a form control (improves accessibility and click targets).
\end{itemize}

\section{3. Lists, Links and Images}
\begin{lstlisting}[language=HTML]
<ol> <!-- ordered list -->
  <li>Football</li>
  <li>Chess</li>
</ol>

<ul> <!-- unordered list -->
  <li>
    <a href="https://example.com">Visit my website</a>
  </li>
</ul>

<img src="https://example.com/img.png" alt="Image">
\end{lstlisting}
\begin{itemize}
    \item \texttt{<ol>} — ordered (numbered) list; \texttt{<ul>} — unordered (bulleted) list; \texttt{<li>} — list item.
    \item \texttt{<a href="...">} creates a hyperlink.
    \item \texttt{<img src="..." alt="...">} embeds an image; \texttt{alt} provides fallback/accessible text.
\end{itemize}

\section{4. Semantic Tags}
\begin{itemize}
    \item \textbf{Non-semantic elements} (\texttt{<div>}, \texttt{<span>}) tell nothing about their content.
    \item \textbf{Semantic elements} (\texttt{<header>}, \texttt{<nav>}, \texttt{<main>}, \texttt{<section>}, \texttt{<article>}, \texttt{<aside>}, \texttt{<footer>}) clearly describe their meaning to both the browser and the developer.
\end{itemize}
\begin{itemize}
    \item \texttt{<header>} — top of the page: logo, title.
    \item \texttt{<nav>} — navigation links.
    \item \texttt{<main>} — the unique content of the page.
    \item \texttt{<section>} — a thematic block with a heading.
    \item \texttt{<article>} — a self-contained piece: a post, a card.
    \item \texttt{<aside>} — side content: menu, related links.
    \item \texttt{<footer>} — contacts, copyright.
\end{itemize}

\section{5. \texttt{<div>} and \texttt{<span>}}
\begin{itemize}
    \item \textbf{\texttt{<div>}} — a block container: takes the full width available, always starts on a new line, groups elements so CSS can style them together.
    \item \textbf{\texttt{<span>}} — an inline container: takes only the width it needs, stays inside the line of text, used to style a few words inside a paragraph.
\end{itemize}
\begin{lstlisting}[language=HTML]
<div class="card"> ... </div>

<p>today is <span class="hot">98</span></p>
\end{lstlisting}

\section{6. Publishing with GitHub Pages}
\begin{enumerate}
    \item Register on GitHub (\url{https://github.com/}).
    \item Create a new repository.
    \item Upload website files — the project's main HTML file must be named \texttt{index.html}.
    \item Enable GitHub Pages in the repository Settings.
    \item Access the published site at the provided URL: \texttt{your-username.github.io/your-repo-name}.
\end{enumerate}

\section{Summary}
Week 2 extends HTML with structured data (tables), user input (forms/inputs/buttons), navigation content (lists, links, images), meaningful page structure (semantic tags), generic grouping containers (\texttt{div}/\texttt{span}), and the practical step of publishing a static site for free via GitHub Pages.

\end{document}
